% ===========================================================
\section{Results}
\label{sec:results}
% ===========================================================

\subsection{Phonon Band Prediction Accuracy}

Table~\ref{tab:main_results} compares the test-set MAE of our
model against several baselines.

\begin{table}[h]
\centering
\caption{Phonon prediction accuracy on the MAX-phase test set.}
\label{tab:main_results}
\begin{tabular}{lcc}
\toprule
Model & MAE (THz) & RMSE (THz) \\
\midrule
Mean predictor (baseline)   & \TODO{X.XX} & \TODO{X.XX} \\
XGBoost (flat features)     & \TODO{X.XX} & \TODO{X.XX} \\
Residual MLP                & \TODO{X.XX} & \TODO{X.XX} \\
Geometric GNN (no physics)  & 0.480       & \TODO{X.XX} \\
\textbf{Ours (Physics-GNN)} & \textbf{\TODO{X.XX}} & \textbf{\TODO{X.XX}} \\
\bottomrule
\end{tabular}
\end{table}

\subsection{Feature Importance Analysis}
\label{sec:feat_results}

Figure~\ref{fig:ablation} shows $\Delta$MAE for the top 10
most important and 5 most harmful atomic features.
The electronegativity \texttt{en\_allen} (+19.66\,cm$^{-1}$)
and Bragg covalent radius \texttt{covalent\_radius\_bragg}
(+17.01\,cm$^{-1}$) are most critical,
consistent with the dominant role of chemical bonding in
determining force constants.
Conversely, \texttt{heat\_of\_formation} ($-$5.14\,cm$^{-1}$)
and redundant radius descriptors degrade accuracy,
likely due to multicollinearity.

% Figure placeholder
\begin{figure}[h]
  \centering
  \includegraphics[width=\columnwidth]{figures/ablation_bar.pdf}
  \caption{Feature ablation: $\Delta$MAE when each atomic descriptor
           is removed. Positive values indicate the feature improves
           prediction accuracy.}
  \label{fig:ablation}
\end{figure}

\subsection{Representative Band Structures}

\TODO{Add figure: predicted vs.~DFT phonon bands for 3--4 representative
MAX phases (one stable, one borderline stable).}

\subsection{Physical Constraint Satisfaction}

The acoustic sum rule residual $\|\sum_j\IFC_{ij}\|$ decreases
from \TODO{X.XX} (unconstrained) to \TODO{X.XX} (with ASR loss term),
confirming that the physics-informed objective improves
physical consistency.
